\section{Method}
\label{sec:method}

% Draft for Ammaar's revision. Week-2 deliverable. Faithful to the
% implementation in src/schematic2netlist/ as of this commit; every
% threshold named here lives in configs/default.yaml (single source of
% truth) — the prose deliberately names the config keys so reviewers
% can map text to code.

Fig.~\ref{fig:pipeline} shows the pipeline: preprocessing, component
detection, wire extraction, net assembly, terminal snapping, netlist
export, and the design-intent repair layer. The design philosophy is
deliberate: a lightweight, locally-runnable system of interpretable
stages, each independently measurable (Sec.~\ref{sec:setup}) and each
switchable to its classical baseline in configuration — the ablations
in Sec.~\ref{sec:results} are config switches, not code forks.

\begin{figure*}[t]
  \centering
  % Pipeline figure (vector art, TikZ) — replaces the legacy bioRender
% flowchart. \input from main via figures/pipeline_figure.
% Two-column width (figure*). Repair layer shaded; config switches dashed.
\begin{tikzpicture}[
  font=\small,
  node distance=4mm and 6mm,
  stage/.style={draw, rounded corners=1pt, align=center,
                minimum height=9mm, inner sep=2.5pt, fill=white},
  data/.style={draw, align=center, minimum height=7mm, inner sep=2.5pt,
               fill=gray!8, font=\footnotesize},
  repair/.style={draw, rounded corners=1pt, align=center,
                 minimum height=9mm, inner sep=2.5pt, fill=orange!15},
  switch/.style={draw, dashed, align=center, inner sep=2pt,
                 font=\footnotesize, fill=white},
  arr/.style={-{Stealth[length=2mm]}, thick},
  lbl/.style={font=\scriptsize\itshape, text=black!60}
]

% ---- main recognition path ----
\node[data] (img) {photo of\\hand-drawn circuit};
\node[stage, right=of img] (pre) {preprocess\\ \scriptsize deskew, shadow,\\ \scriptsize binarize, crop};
\node[stage, right=of pre] (det) {detect\\ \scriptsize YOLOv8s, 17 classes\\ \scriptsize incl.\ Wire Crossover};
\node[stage, right=of det] (wires) {wire extraction\\ \scriptsize ink mask, masking,\\ \scriptsize hole stitching};
\node[stage, right=of wires] (nodes) {net assembly\\ \scriptsize crossover-aware\\ \scriptsize connected comp.};
\node[stage, right=of nodes] (snap) {terminal\\snapping\\ \scriptsize boundary};

\draw[arr] (img) -- (pre);
\draw[arr] (pre) -- (det);
\draw[arr] (det) -- (wires);
\draw[arr] (wires) -- (nodes);
\draw[arr] (nodes) -- (snap);

% detector boxes feed masking and crossover handling
\draw[arr, black!50] (det.south) to[out=-60,in=-120]
  node[lbl, below] {component/text masks \; \textbullet\; crossover boxes} (nodes.south);

% ---- second row: export + repair + outputs ----
\node[stage, below=9mm of nodes] (netlist) {netlist export\\ \scriptsize role-based SPICE,\\ \scriptsize model cards};
\node[repair, left=of netlist] (rep) {ERC + minimal repair\\ \scriptsize gauge vs.\ assumption,\\ \scriptsize topology untouched};
\node[data, left=of rep] (ledger) {assumption\\ledger (JSON)};
\node[data, below=6mm of rep] (spice) {SPICE netlist\\ \scriptsize ngspice DC check};

\draw[arr] (snap.south) to[out=-90,in=20] (netlist.east);
\draw[arr] (netlist) -- (rep);
\draw[arr] (rep) -- (ledger);
\draw[arr] (rep) -- (spice);

% ---- ablation switches ----
\node[switch, above=2.5mm of wires] (sw1) {\texttt{wires.method}\\ink $\leftrightarrow$ canny};
\node[switch, above=2.5mm of nodes] (sw2) {\texttt{handle\_crossovers}\\on $\leftrightarrow$ off};
\node[switch, above=2.5mm of snap] (sw3) {\texttt{snapping.strategy}\\boundary / directional / uniform};
\draw[dashed, black!60] (sw1) -- (wires);
\draw[dashed, black!60] (sw2) -- (nodes);
\draw[dashed, black!60] (sw3) -- (snap);

\end{tikzpicture}

  \caption{The pipeline. Solid stages form the recognition path;
  the shaded stage is the C5 design-intent repair layer, which may
  add explicit, logged constraints but provably never modifies the
  recovered topology. Configuration switches (dashed) select the
  classical baselines used in the C2 ablation.}
  \label{fig:pipeline}
\end{figure*}

\subsection{Component Detection}
\label{sec:method-detect}
A YOLOv8s \cite{yolov8} detector is trained on the Digitize-HCD train
split at $640$\,px over the 17-class vocabulary (three seeds;
Sec.~\ref{sec:setup}). Detection includes the \emph{Wire Crossover}
class — a drawing convention marking where two wires cross without
connecting — which most prior work discards but which our net assembly
exploits (Sec.~\ref{sec:method-nodes}). Detections are cached
per-image, making every downstream experiment reproducible from the
same detection set.

\subsection{Wire Extraction}
\label{sec:method-wires}
The wire mask is computed from binarized \emph{ink} rather than Canny
edges: hand-drawn strokes are solid, so edge detection yields hollow
double-contours that fragment under morphology (the classical
\texttt{canny} path remains available as the ablation baseline,
\texttt{wires.method}). Regions belonging to detected non-wire classes
(component bodies, padded by 8\,px) and detected text are masked out —
except \emph{Wire Crossover} boxes, which would sever the very wires
that cross there. Anisotropic along-stroke bridging closes pen gaps up
to 9\,px, and a blob filter keeps thin strokes by extent rather than
area alone.

\subsubsection{Stitching Self-Inflicted Mask Holes}
Masking components and text out of the ink necessarily cuts the wires
that touch them. Measurement on the benchmark showed this was the
\emph{dominant} connectivity failure: supply rails shattered into 4--6
islands with 20--45\,px gaps at exactly the masked rectangles. The
stitching stage reconnects wire islands across those known holes under
three conservative guards: (i) the bridge must be collinear with the
local stroke direction at \emph{both} endpoints (within
$35^\circ$), (ii) the gap must be explained by a masked rectangle it
traverses, and (iii) the bridge must not pass through a third island.
Two passes handle chained holes (a text box adjacent to a component
pad). Crucially, stitching never bridges across a component
\emph{body} — a component's two leads are different nets by
definition.

\subsection{Net Assembly: Crossover-Aware Connected Components}
\label{sec:method-nodes}
Classically, each 8-connected component of the wire mask is one
electrical net. This has a structural ceiling: two wires that
\emph{cross without connecting} form a single pixel-connected region
and are merged into one net — no threshold can fix this. We use the
detector's \emph{Wire Crossover} boxes to break the ceiling: inside
each detected crossover, the mask is notched to separate the crossing
strokes, and the four wire arms entering the box are re-linked
pairwise (opposite arms belong to the same net; adjacent arms do not).
The result is a net partition that respects drawn crossing semantics.
The switch (\texttt{nodes.handle\_crossovers}) is the core of the C2
ablation.

\subsection{Terminal Snapping}
\label{sec:method-snap}
Each component must attach its terminals to nets. The
\emph{boundary-crossing} strategy walks the component's bounding-box
boundary and collects the wire strokes that cross it, using the
class's expected terminal count (two-terminal passives and sources,
three-terminal transistors, five-terminal op-amps) to select terminal
sites. Two legacy strategies — \emph{directional} (orientation-guessed
snap windows) and \emph{uniform} (stepwise bbox expansion) — are kept
as configuration baselines for the snapping ablation.

\subsection{Port-Template Pin Identity and Polarity (C3)}
\label{sec:method-ports}
Finding \emph{where} a component connects is not the same as knowing
\emph{which pin} connects there. Under the strategies above, terminals
are filled in whatever order the geometry search encounters them, so a
diode's anode and cathode — or a MOSFET's drain, gate and source — are
assigned arbitrarily, and the emitted netlist cannot be trusted for any
directional device.

Digitize-HCD supplies the missing supervision, in an annotation
modality no published work exploits: per class, thousands of component
crops with pixel port coordinates and, for directional parts,
\emph{port names} (Anode/Cathode, Positive/Negative, Drain/Gate/Source,
In$+$/In$-$/Out). We distill these into per-class \emph{port
templates}: the median normalized position of each named port, in
canonical order.

The modeling choice that makes such templates useful is the binning.
Hand-drawn symbols appear in any orientation, and splitting merely into
horizontal and vertical merges mirror images — a diode pointing left
and one pointing right are both ``horizontal'' — which places every
port near the middle of the crop. We therefore bin by \emph{pose}: the
angle of the vector between the first and last port in canonical order,
quantized into eight sectors, derived from the port constellation
itself and never from image content. At inference, boundary crossings
are located as above, then Hungarian-matched against each candidate
pose's predicted pin sites; the best-fitting pose supplies both the
assignment and the pin identities. When no pose fits, the component
falls back to boundary order, so pin identity is a bonus and never a
connectivity regression. Because the template port order is
deliberately SPICE argument order, the writer's positional emission
(\texttt{M d g s}, \texttt{Q c b e}, \texttt{D anode cathode}) then
carries its intended meaning.

\subsection{Netlist Export}
\label{sec:method-netlist}
Nets are named (ground is net \texttt{0}, chosen by GND-symbol
snapping with a configurable fallback policy), and each component is
emitted as a SPICE element according to its class \emph{role};
referenced device models get \texttt{.model} cards so diode/MOSFET/BJT
netlists are parseable. Component \emph{values} are not read from the
drawing (explicit future work); placeholders are assigned — and,
importantly for C5, that assignment is \emph{logged as an assumption},
not silent (Sec.~\ref{sec:method-repair}).

\subsection{Design-Intent Repair with an Assumption Ledger (C5)}
\label{sec:method-repair}
A topologically correct netlist can still fail simulation: no ground
reference, a subnet with no DC path to ground, an ideal-inductor DC
short, an unset source polarity. These are \emph{underspecifications}
of the drawing, not recognition errors, and prior systems either fail
silently or fix silently. Our repair layer makes this step explicit,
minimal, and auditable.

\textbf{Diagnosis (ERC).} Rule checks on the recovered graph mirror
the reasons simulators reject circuits, using DC-conduction reasoning
(capacitors and current sources open at DC; inductors and sources
conduct; transistor channels conduct, gates/bases do not — a
documented approximation adequate for diagnosis).

\textbf{Repair with a gauge/assumption taxonomy.} Each issue maps to a
strategy that emits a \emph{ledger entry} classified as either
\begin{itemize}
  \item \emph{gauge} — behavior-invariant or determined by the drawing
  (current-source reference direction, terminal order of a symmetric
  passive, ground selection when a GND symbol exists): safe to infer,
  still logged; or
  \item \emph{assumption} — behavior-changing (choosing a reference
  node when no ground is drawn, shunting a floating subnet,
  placeholder component values): applied only if needed, flagged with
  alternatives and a confidence, and counted against a minimality
  budget (\texttt{repair.max\_assumptions}).
\end{itemize}
Strategies prefer standard SPICE convergence aids \cite{ngspice} over
structural edits, and the layer applies the fewest, lowest-impact
fixes that yield DC solvability.

\textbf{The integrity rule.} Repair \emph{never} changes the recovered
topology — it only adds explicit constraints. This is enforced by
construction, unit-tested, and verified experimentally: topology
metrics are identical with repair on and off
(Sec.~\ref{sec:results}). The per-circuit ledger (JSON schema v1,
released with the benchmark) records issue, category, action,
alternatives, confidence, and evidence for every entry — the complete,
machine-readable answer to ``what did the system assume about my
circuit?''
